% Generated by Listas/make_statements.py from pratica-lecture-07.md, item 18. Do not edit.
\textbf{Source.} Tadelis (2013), \emph{Game Theory: An Introduction}, Exercise 12.8 ``Jury Voting'', pp.~268--269 {[}PDF pp.~287--288{]}. Reproduced verbatim. The game it refers to is Tadelis §12.4 (pp.~261--264), summarized in Notas/L07/L07-09-common-values-curses §1. \textbf{Answer in the source:} no (none in the book; Princeton UP's website has about half, not checked).



\begin{quote}

\textbf{Jury Voting:} Consider the jury voting game in Section 12.4.



\begin{enumerate}

\def\labelenumi{\alph{enumi}.}

\item

  Show that both players choosing \(CC\) is a Bayesian Nash equilibrium of the two-player game.

\item

  Would it be better for the two players if only one of them decided on the fate of the defendant? Why or why not?

\item

  Consider a three-player game in which the defendant is convicted if at least two players vote to convict him, and acquitted otherwise (i.e., majority voting determines the defendant's fate). Assume that the informational structure is the same as in the two-player game. Are there conditions in which voting according to one's signal (\(C\) if \(\theta_i = \theta_G\) and \(A\) if \(\theta_i = \theta_I\)) is a Bayesian Nash equilibrium?

\end{enumerate}

\end{quote}



\textbf{Note on the book (item 18).} Part (a) is true only for \(p \le q\) (the signal precision does not exceed the

prior): against a juror who always convicts, a juror who sees \(\theta_I\) is pivotal and acquits when \(p > q\).

The book (p.~264) states \(CC\) as an equilibrium without the condition. See the solutions PDF, item 18.



\begin{center}\rule{0.5\linewidth}{0.5pt}\end{center}

